% =============================================================================
% БИЛЕТ 21
% =============================================================================

\ticket{21}{}

\question{Касательное пространство к гладкому многообразию. Корректность определения}

\begin{defn}[Гладкое многообразие]
  Множество $M \subset \R^n$ называется \emph{гладким многообразием}, если для каждой точки $p \in M$ существует окрестность $U(p)$ такая, что $M \cap U(p)$ --- простое гладкое многообразие.
\end{defn}

\begin{defn}[Касательная к многообразию]
  Пусть $M \subset \R^n$ --- гладкое многообразие, $p \in M$. Выберем такую точку $p \in M$ и функцию $\gamma : \in \mathbb{C}^1([a, b], \R^n)$ такую, что $\gamma(t) \in M$ для всех $t \in [a, b]$ и $\gamma(t_0) = p$ для некоторого $t_0 \in (a, b)$. Тогда вектор $\gamma'(t_0)$ называется \emph{касательным вектором} к многообразию $M$ в точке $p$.
\end{defn}
\begin{defn}[Касательное пространство]
  Множество всех касательных векторов к многообразию $M$ в точке $p$ называется \emph{касательным пространством} к многообразию $M$ в точке $p$ и обозначается $T_p M$.
\end{defn}
\question{Признаки Дирихле и Абеля сходимости несобственных интегралов}

% Ответ...

